% Inactive fragment: bounded domination in the invariant Hessian mass regime.
% The global chapter setting remains compact X, Hermitian omega, 2 <= m < n.
\section{Bounded domination under the power identities}
\label{draft629:bounded-domination}

The positive background term in a convex rescaling makes a strict
comparison set visible to volume. Under the power identities of
\cref{hessian:global-comparison-characterization}, the global comparison
principle then yields domination for bounded potentials. Canonical
representatives turn the resulting almost everywhere order into
pointwise order.

\begin{lemma}\label{draft629:canonical-order}
If bounded admissible $u,v$ satisfy $u\leq v$ almost everywhere
with respect to $\omega^n$, then $u\leq v$ everywhere.
\end{lemma}
\begin{proof}
By \cref{hessian:regularized-maximum}, $w=\max\{u,v\}$ is bounded
admissible with its indicated canonical representative. Since
$w=v$ almost everywhere, they represent the same distribution.
The canonical subharmonic representative used in
\eqref{eq:hessian-bounded-upper-tests} is determined by that
distribution. Thus $w=v$ pointwise, which is the assertion.
\end{proof}

\begin{lemma}\label{draft629:convex-background-term}
For bounded admissible $v$, $0\leq\alpha\leq1$ and a real constant
$b$, the function $v_\alpha=\alpha v+b$ is bounded admissible and
\[
 H_m(v_\alpha)\geq\alpha^m H_m(v)+(1-\alpha)^m\omega^n.
\]
\end{lemma}
\begin{proof}
Take smooth decreasing uniformly bounded admissible $v_j\downarrow v$.
Cone convexity makes $\alpha v_j+b$ admissible because
\[
 \omega+\ddc(\alpha v_j+b)
        =\alpha(\omega+\ddc v_j)+(1-\alpha)\omega.
\]
For $\alpha>0$, an upper test $q$ for $\alpha v+b$ gives the
upper test $(q-b)/\alpha$ for $v$. The same convex combination
of forms therefore proves admissibility of the limit. For
$\alpha=0$ the potential is constant. The pointwise decreasing
limit has the displayed canonical affine values.
Expand the smooth top product. Every mixed term is positive by
polarized G\aa rding positivity. Keeping the two extreme terms gives
\[
 H_m(\alpha v_j+b)\geq\alpha^m H_m(v_j)
                                    +(1-\alpha)^m\omega^n.
\]
Joint decreasing convergence from
\cref{hessian:joint-decreasing-products} passes this inequality
against each nonnegative continuous test to the asserted inequality
of positive Radon measures.
\end{proof}

\begin{theorem}\label{draft629:strict-bounded-domination}
Assume $\ddc(\omega^{n-q})=0$ for $1\leq q\leq m$. Let $u,v$
be bounded admissible. If, for some $0\leq c<1$,
\[
 \mathbf1_{\{u<v\}}H_m(u)
                   \leq c\mathbf1_{\{u<v\}}H_m(v),
\]
then $u\geq v$ everywhere. In particular, this conclusion holds
if $H_m(u)(\{u<v\})=0$.
\end{theorem}
\begin{proof}
Choose $B>\|v\|_\infty$ and, for $0<\alpha<1$ with $\alpha^m>c$,
put $v_\alpha=\alpha v-(1-\alpha)B$. Then $v_\alpha\leq v$ and
$E_\alpha=\{u<v_\alpha\}\subset\{u<v\}$.
By \cref{hessian:global-comparison-characterization} and the
assumed restriction inequality,
\[
 \int_{E_\alpha}H_m(v_\alpha)
        \leq\int_{E_\alpha}H_m(u)
        \leq c\int_{E_\alpha}H_m(v).
\]
On the other hand, \cref{draft629:convex-background-term} gives
\[
 (\alpha^m-c)\int_{E_\alpha}H_m(v)
             +(1-\alpha)^m\int_{E_\alpha}\omega^n\leq0.
\]
Both summands are nonnegative, and the coefficient of the volume
is positive. Thus $E_\alpha$ has zero volume. Applying
\cref{draft629:canonical-order} gives $u\geq v_\alpha$ pointwise.
Let $\alpha$ increase to one. The boundedness of $v$ gives
$v_\alpha\to v$ uniformly, so $u\geq v$ everywhere. The final
assertion is the case $c=0$.
\end{proof}

\begin{corollary}\label{draft629:proper-bounded-uniqueness}
Assume $\ddc(\omega^{n-q})=0$ for $1\leq q\leq m$ and fix $a>0$.
If bounded admissible $u,v$ satisfy
\[
 e^{-av}H_m(v)\geq e^{-au}H_m(u),
\]
then $u\geq v$. Consequently, for any fixed finite positive Radon
measure $\nu$, the equation $H_m(u)=e^{au}\nu$ has at most one
bounded admissible solution.
\end{corollary}
\begin{proof}
For $\delta>0$, put $E_\delta=\{u<v-\delta\}$. Multiplying the
assumed measure inequality by the positive bounded Borel function
$e^{au}$ gives on $E_\delta$
\[
 H_m(u)\leq e^{a(u-v)}H_m(v)
                 \leq e^{-a\delta}H_m(v-\delta).
\]
The last identity uses invariance under constants. Since
$e^{-a\delta}<1$, \cref{draft629:strict-bounded-domination} applied
to $u,v-\delta$ gives $u\geq v-\delta$ everywhere. Letting
$\delta\downarrow0$ gives $u\geq v$. Two solutions of the stated
equation have equal weighted measures $e^{-au}H_m(u)=\nu$.
Applying the order assertion in both directions proves uniqueness.
\end{proof}
